\documentclass[10pt]{article}
\usepackage[T1]{fontenc}
\usepackage[utf8]{inputenc}
\usepackage[french]{babel}
\usepackage{amsmath,amssymb,mathrsfs,mathpazo}
\usepackage{geometry}
\geometry{a4paper,margin=2cm}
\usepackage{graphicx,array,tabularx,tabularray,colortbl,multirow,multicol,booktabs,enumitem,pifont,ulem,fancybox,diagbox,cancel}
\usepackage{pstricks,pstricks-add,pst-plot,pst-tree,pst-node,pst-text,pst-func,pst-math,pst-bspline,pst-3dplot,pst-eucl,pst-solides3d}
\let\pscalculate\undefined
\usepackage{pst-bezier}
\usepackage{tikz,pgfplots,tkz-tab}
\usetikzlibrary{arrows,arrows.meta,patterns,calc,positioning,trees,decorations.markings,shapes.geometric,shapes.misc,shapes.arrows,intersections,decorations.pathmorphing,decorations.pathreplacing,backgrounds,shadows,babel}
\pgfplotsset{compat=1.18}
\usepgfplotslibrary{fillbetween}
\usepackage[np]{numprint}
\usepackage{listings,makecell,pas-tableur,twoopt}
\let\barre\undefined
\usepackage[pstricks]{tablvar}
\definecolor{darkgreen}{RGB}{0,100,0}
\providecommand{\red}{\color{red}}
\providecommand{\blue}{\color{blue}}
\providecommand{\green}{\color{green}}
\providecommand{\black}{\color{black}}
\providecommand{\white}{\color{white}}
\providecommand{\gray}{\color{gray}}
\providecommand{\cyan}{\color{cyan}}
\providecommand{\magenta}{\color{magenta}}
\providecommand{\yellow}{\color{yellow}}
\providecommand{\violet}{\color{violet}}
\providecommand{\orange}{\color{orange}}
\providecommand{\euro}{\texteuro}
\providecommand{\EUR}{\texteuro}
\providecommand{\ieme}{\textsuperscript{e}}
\providecommand{\degres}{\ensuremath{^\circ}}
\providecommand{\degree}{\ensuremath{^\circ}}
\providecommand{\vv}[1]{\overrightarrow{#1}}
\providecommand{\Coord}[1]{#1}
\providecommand{\decosix}{}
\providecommand{\starredbullet}{\ensuremath{\bullet}}
\pagestyle{empty}
\setlength{\parindent}{0pt}
\usepackage{adjustbox}\DeclareUnicodeCharacter{2236}{:}\DeclareUnicodeCharacter{2190}{\ensuremath{\leftarrow}}
\providecommand{\R}{\mathbb{R}}
\providecommand{\N}{\mathbb{N}}
\providecommand{\Z}{\mathbb{Z}}
\providecommand{\Q}{\mathbb{Q}}
\providecommand{\D}{\mathbb{D}}
\providecommand{\C}{\mathbb{C}}
\providecommand{\e}{{\mathrm{e}}}
\providecommand{\vect}[1]{\overrightarrow{#1}}
\providecommand{\vectt}[1]{\overrightarrow{\mathrm{#1}}}
\providecommand{\vvt}[1]{\overrightarrow{\mathrm{#1}}}
\providecommand{\np}[2][]{{#2~#1}}
\providecommand{\vv}[1]{{\overrightarrow{#1}}}
\renewcommand{\d}{{\mathrm{d}}}
\providecommand{\degres}{{^\circ}}
\providecommand{\degre}{{^\circ}}
\providecommand{\blue}{\color{blue}}
\providecommand{\red}{\color{red}}
\providecommand{\enskip}{\quad}
\providecommand{\llbracket}{\lbrack\!\lbrack}
\providecommand{\rrbracket}{\rbrack\!\rbrack}
\providecommand{\euro}{{\text{€}}}
\providecommand{\no}{{\text{n°}}}
\providecommand{\interval}[2]{{[#1;#2]}}
\providecommand{\intervFO}[2]{{[#1;#2[}}
\providecommand{\dotfill}{\ldots}

\providecommand{\ds}{\displaystyle}
\providecommand{\textding}[1]{#1}
\providecommand{\Oij}{\(\mathrm{O};\vec{\imath},\vec{\jmath}\)}
\providecommand{\Oijk}{\(\mathrm{O};\vec{\imath},\vec{\jmath},\vec{k}\)}
\providecommand{\Ouv}{\(\mathrm{O};\vec{u},\vec{v}\)}
\providecommand{\intervFF}[2]{\ensuremath{[#1;#2]}}
\renewcommand{\euro}{\ifmmode\text{\texteuro}\else\texteuro\fi}
\DeclareUnicodeCharacter{00B0}{\ifmmode^{\circ}\else\textdegree\fi}
\DeclareUnicodeCharacter{20AC}{\ifmmode\text{\texteuro}\else\texteuro\fi}
% Énoncés : Éducation nationale. Transcription et corrigés : APMEP, Denis Vergès et contributeurs. Découpage et classement : CoopMaths, Sébastien Lozano, Éric Elter et contributeurs. Adaptation de la présentation : Maths Explorer.
% CC BY-SA : https://coopmaths.fr/www/about/
% https://coopmaths.fr/alea/static/eam/2026/tex/eam_specialite_2026_06_metropole_1_cor.tex
\begin{document}
\noindent\textit{APMEP / CoopMaths. CC BY-SA. Présentation adaptée par Maths Explorer.}\par\medskip
Un loueur de bicyclettes propose deux types de bicyclettes : des bicyclettes traditionnelles et des bicyclettes électriques. Il incite ses clients à prendre une assurance. On dispose des informations suivantes.
\begin{itemize}[label=$\bullet$]
	\item 60\% des clients ont loué une bicyclette traditionnelle, les autres ont loué une bicyclette électrique.
	\item Parmi ceux qui ont loué une bicyclette traditionnelle, 25\% ont pris une assurance.
	\item 20\% de l'ensemble des clients ont pris une assurance.
\end{itemize}

On choisit un client au hasard et on note les événements :

\textit{T} : \og{}le client a loué une bicyclette traditionnelle\fg{}.

\textit{A} : \og{}le client a pris une assurance\fg{}.

\begin{enumerate}
	\item Recopier l'arbre pondéré ci-dessous et compléter les pointillés.


\begin{center}
			\begin{tikzpicture}[xscale=1,yscale=1,rotate=0,every node/.style={scale=1,rotate=0}]
			\node(n0)at(0,0){};
			\node(n1_1)at(3,1){$T$};
			\node(n1_2)at(3,-1){$\overline{T}$};
			\node(n2_1)at(6,1.5){$A$};
			\node(n2_2)at(6,0.5){$\overline{A}$};
			\node(n2_3)at(6,-0.5){$A$};
			\node(n2_4)at(6,-1.5){$\overline{A}$};
			\draw[blue](n0)--(n1_1)node[pos=0.6,fill=white]{\np{0.6}};
			\draw[blue](n0)--(n1_2)node[pos=0.4,fill=white]{\np{0.4}};
			\draw[blue](n2_1)--(n1_1)node[pos=0.6,fill=white]{\np{0.25}};
			\draw[blue](n2_2)--(n1_1)node[pos=0.4,fill=white]{\np{0.75}};
			\draw[blue](n2_3)--(n1_2);
			\draw[blue](n2_4)--(n1_2);

		\end{tikzpicture}

\end{center}
		\item Donner, par simple lecture de l'énoncé, la probabilité de l'événement A.


		 \textcolor{red}{\textsc{Correction}}


		20\% des clients ont pris l'assurance, donc \mbox{$P(A)=\np{0.2}$}
		\item Montrer que la probabilité que le client ait loué une bicyclette traditionnelle et qu'il ait pris une assurance est égale à \np{0.15}.

		 \textcolor{red}{\textsc{Correction}}


		\mbox{$P(T \cap A)=P(T)\times P_T(A)=\np{0.6}\times\np{0.25}=\np{0.15}$}.

		\item En déduire que la probabilité \mbox{$P(\overline{T}\cap A)$} est égale à \np{0.05}.

		 \textcolor{red}{\textsc{Correction}}


		Loi des probabilités totales : \quad \mbox{$P(A)=P(T\cap A)+P(\overline{T}\cap A)$};

		donc \mbox{$P(\overline{T}\cap A)=P(A)-P(T\cap A)=\np{0.2}-\np{0.15}=\np{0.05}$}

		la probabilité qu'un client ait loué une bicyclette électrique et pris une assurance est \np{0.05}.

		\item Déterminer la probabilité que le client ait pris une assurance sachant qu'il a loué une bicyclette électrique. On donnera le résultat sous la forme d'une fraction irréductible.

			 \textcolor{red}{\textsc{Correction}}


		Puisque \quad \mbox{$P(\overline{T}\cap A)=P(\overline{T})\times P_{\overline{T}}(A)$},

		On en déduit :\quad \mbox{$P_{\overline{T}}(A)=\dfrac{P(\overline{T} \cap A)}{P(\overline{T})} = \dfrac{\np{0.05}}{\np{0.4}}=\dfrac18$}

		la probabilité qu'un client ayant loué une bicyclette électrique ait pris une assurance est \mbox{$\dfrac18=\np{0.125}$}

\end{enumerate}

\end{document}
